\documentclass[12pt]{article}
\usepackage{graphicx}
\usepackage[spanish]{babel}
\usepackage[a4paper]{geometry}
\usepackage{hyperref}
\usepackage[none]{hyphenat}
\usepackage{listings}
\usepackage{color}
\hypersetup{
    colorlinks=true,
    linkcolor=black,
    filecolor=magenta,      
    urlcolor=cyan,
    pdfpagemode=FullScreen,
    }
\definecolor{deepblue}{rgb}{0,0,0.5}
\definecolor{deepred}{rgb}{0.6,0,0}
\definecolor{deepgreen}{rgb}{0,0.5,0}
\lstdefinestyle{pythonstyle}{
    language=Python,
    basicstyle=\ttfamily\small,
    keywordstyle=\color{deepblue},
    stringstyle=\color{deepgreen},
    emphstyle=\ttb\color{deepred},
    showstringspaces=false,
    breaklines=true,
    frame=single
}
\pagestyle{headings}
\title{Pinocho Univalle}
\begin{document}
\begin{titlepage}
\centering
{\includegraphics[width=0.2\textwidth]{images/logo.png}\par}
\vspace{1cm}
{\bfseries\LARGE Universidad del Valle \par}
\vspace{1cm}
{\scshape\Large Ingeniería de Sistemas \par}
\vfill
{\scshape\Huge Pinocho Univalle \par}
\vfill
{\itshape\Large Introducción a la Inteligencia Artificial \par}
\vfill
{\Large Autores: \par}
{\Large Andrés Felipe Ruíz Buriticá \par}
{\Large Kevin Steven Victoria Ospina \par}
{\Large Juan Sebastián González Camacho \par}
\vfill
{\Large \today \par}
\end{titlepage}
\tableofcontents
\clearpage
\section{Introducción}
Un agente inteligente, es una entidad capaz de percibir su entorno, procesar tales percepciones y responder o actuar en su entorno de manera racional, es decir, de manera correcta y tendiendo a maximizar un resultado esperado. Es capaz de percibir su medio ambiente con la ayuda de sensores y actuar en ese medio utilizando actuadores.
\newline \\
En este proyecto, implementaremos algunos de los algoritmos de búsqueda no informada, aplicados a los agentes inteligentes, que aprendimos en lo que lleva del curso de Introducción a la Inteligencia Artificial.
\clearpage
\section{Objetivo}
Este proyecto tiene como objetivo aplicar los conceptos vistos en el transcurso del curso desarrollando un agente inteligente para Pinocho que le ayude a encontrar a Gepetto. Para lograrlo debemos tener en cuenta los siguientes puntos:
\begin{itemize}
    \item El agente se mueve a los espacios vacíos o hacia Gepetto con costo 1.
    \item Si el agente pasa por los cigarrillos, cuesta 2.
    \item Si el agente pasa por algún zorro, cuesta 3.
    \item La matriz debe leerse desde un archivo de texto.
    \item Se debe implementar la técnica de amplitud (en este caso aplicando una variante en la que se recorra en zig zag), costo uniforme y profundidad iterativa.
    \item Debe al final definir con buenos criterios, cuál de las tres estrategias fue mejor y por qué.
\end{itemize}
\clearpage
\section{Desarrollo}
\begin{enumerate}
    \item Creamos una clase \textit{Maze} para almacenar el tablero de juego, que se puede leer desde una matriz numérica dentro de un archivo de texto. Para esto, hacemos uso de la librería \textit{Numpy}, que nos permite hacer la lectura del archivo de texto y representarlo como un arreglo. Por último, se define el punto de partida (la posición de Pinocho) y el punto de destino (la posición de Gepetto).\\
\begin{lstlisting}[style=pythonstyle]
import numpy as np
from constant import Constant

class Maze(Constant):
    """
    Class that represents a maze.
    """

    def __init__(self, filename:str, matrix = None):
        """
        Initializes the class instance with a numeric
        matrix from a file.

        Args:
            filename (str): path to the file with the
            numeric matrix representing the maze.
        """
        if not matrix is None:
            self.maze = np.array(matrix)
        else:
            # Load maze from file.
            self.maze = np.loadtxt(filename, dtype=int)
            
        self.matrix = self.maze.tolist()

        # Define the start and goal positions.
        self.start = tuple(
            np.argwhere(self.maze == self.PINOCCHIO)[0])
        self.goal = tuple(
            np.argwhere(self.maze == self.GEPETTO)[0])
\end{lstlisting}
\end{enumerate}
\clearpage
\section{Conclusión}
\clearpage
\end{document}